\section{Rooted trees and cancellation at the positive singularity}\label{sec:positive}
\subsection{The singular logarithm of an integer polylogarithm}
For an integer $r\geq1$, the local expansion on a fixed slit neighborhood is
\begin{align}\label{eq:polylog}
\Li_r(\ee^{-rs})={}&
\frac{(-rs)^{r-1}}{(r-1)!}\bigl(H_{r-1}-\log(rs)\bigr)\\
&+\sum_{\substack{k\geq0\\k\ne r-1}}
 \zeta(r-k)\frac{(-rs)^k}{k!},\notag
\end{align}
where $H_0=0$ and $H_m=\sum_{h=1}^m h^{-1}$. For each of the finitely many $r\leq L$, the analytic part is convergent sufficiently near $s=0$. With a fixed choice of $\log s$, \eqref{eq:polylog} is valid on both sides of the short boundary arc under consideration. The singularity comes from the one displayed logarithm.

Multiplying by the coefficient $(-1)^{r+1}\ee^{-rs}/r$ from $H_L$ gives the logarithmic part
\[
-\frac{r^{r-1}}{r!}\,s^{r-1}\ee^{-rs}\log s.
\]
Both sign factors $(-1)^{r+1}$ and $(-1)^{r-1}$ have canceled; the remaining minus sign is from $-\log(rs)$. Hence, at $z=\ee^{-s}$, the logarithmic coefficient of $H_L$ is $-B_{1,L}(s)$, where
\begin{equation}\label{eq:B1}
B_{1,L}(s)=\sum_{r=1}^{L}\frac{r^{r-1}}{r!}\,
 s^{r-1}\ee^{-rs}.
\end{equation}

\subsection{The exact tree identity}
Let
\begin{equation}\label{eq:tree}
T(w)=\sum_{r\geq1}\frac{r^{r-1}}{r!}w^r.
\end{equation}
Lagrange inversion gives $T(w)=w\ee^{T(w)}$ for the analytic branch with $T(0)=0$. Equivalently, $T(w)\ee^{-T(w)}=w$. Since $s\mapsto s\ee^{-s}$ has derivative $1$ at zero, the inverse function theorem implies
\begin{equation}\label{eq:treeidentity}
T(s\ee^{-s})=s
\end{equation}
for $s$ in a neighborhood of zero. Dividing its power series by $s$, with the removable value at $s=0$, yields
\begin{equation}\label{eq:exactcancel}
\sum_{r\geq1}\frac{r^{r-1}}{r!}\,
 s^{r-1}\ee^{-rs}=1.
\end{equation}
Every fixed Taylor coefficient receives contributions from finitely many $r$. Consequently
\begin{equation}\label{eq:D}
B_{1,L}(s)=1+O(s^L),\qquad D_L(s):=1-B_{1,L}(s)=O(s^L).
\end{equation}
No interchange of a nonconvergent boundary logarithmic series is used: the finite factorization is already exact, and \eqref{eq:exactcancel} is an analytic identity near zero.

\begin{proposition}[Only a pole remains at the positive root]\label{prop:positive}
For every finite $L$, there is a function $U_L$ analytic and nonzero near zero such that
\begin{equation}\label{eq:positiveform}
P_L(\ee^{-s})=s^{-1}U_L(s)
 \exp\bigl(D_L(s)\log s\bigr).
\end{equation}
The factor $Q_L(\ee^{-s})$ has the differentiated finite Taylor expansions of Lemma~\ref{lem:taylor}. It follows that, to every prescribed finite order, the only nonanalytic coefficient-producing term at $z=1$ is a constant multiple of $(1-z)^{-1}$.
\end{proposition}
\begin{proof}
The terms of $\log P_L(\ee^{-s})$ other than $-B_{1,L}(s)\log s$ are analytic near zero. This includes $\log(1+\ee^{-2s})$, whose argument is nonzero there. Exponentiating the analytic part produces $U_L$. Since $-B_{1,L}=-1+D_L$, equation \eqref{eq:positiveform} follows.

Now fix the desired finite order, then choose $L$ larger. The logarithmic exponential differs from $1$ first at order $s^L\log s$, so after multiplication by $s^{-1}$ its first possible logarithm is at order $s^{L-1}\log s$. Before that order, multiplication by a sufficiently deep Taylor polynomial of $Q_L$ produces a Laurent polynomial with one simple pole and otherwise nonnegative integer powers. The change of variable $s=-\log z$ is analytic and invertible near $z=1$, and $s^{-1}-(1-z)^{-1}$ is analytic there. Thus the sole singular contribution at the selected orders is the simple pole.
\end{proof}

The conclusion is a finite-jet assertion for each requested differentiability order. It is not a convergent local Laurent expansion of $F$ in a full neighborhood of $1$. Dense singularities on the unit circle will rule out that stronger claim.

\subsection{The pole coefficient from the product}
For real $0<z<1$,
\begin{align}\label{eq:residue}
\log\bigl((1-z)F(z)\bigr)
={}&(1-z)\log(1-z)\\
&+\sum_{j\geq1}\left\{\log\left(1+\frac{z^{j+1}}j\right)
-\frac{z^{j+1}}j\right\}.\notag
\end{align}
For $j\geq2$ the summand is uniformly bounded in absolute value by a fixed multiple of $j^{-2}$; the first summand is handled separately. Dominated convergence gives
\begin{align*}
\lim_{z\uparrow1}\log\bigl((1-z)F(z)\bigr)
&=\sum_{j\geq1}\left(\log(1+1/j)-1/j\right)\\
&=\lim_{N\to\infty}\bigl(\log(N+1)-H_N\bigr)=-\gamma.
\end{align*}
Therefore
\begin{equation}\label{eq:resvalue}
\lim_{z\uparrow1}(1-z)F(z)=\ee^{-\gamma}=A.
\end{equation}
The finite-jet form identifies this radial limit with the coefficient of the local pole. This is a direct residue computation, rather than a Tauberian inference from the coefficient limit. Since $[z^m](1-z)^{-1}=1$, the positive root contributes exactly $A$ to every coefficient in the final subtraction.
